\documentclass[11pt,UTF8,fontset=fandol]{ctexart}
\usepackage{amsmath,amssymb,amsfonts}
\usepackage{graphicx,float,caption,subcaption}
\usepackage{booktabs,array,multirow,makecell}
\usepackage{geometry,enumitem,xcolor}
\usepackage[hidelinks]{hyperref}
\usepackage{url}
\geometry{margin=2.3cm}
\graphicspath{{/Users/yye/project/论文/assets/instructgraph/}}
\captionsetup{font=small,labelfont=bf}
\setlength{\parskip}{0.4em}
\hypersetup{pdftitle={InstructGraph 中文精读导读},pdfauthor={Wang et al.}}

% 小白知识卡片：上下加线 + 粗体标签，稳健不挑内容
\newenvironment{tip}[1]{%
  \par\smallskip\noindent\rule{\linewidth}{0.4pt}\par\nobreak
  \noindent\textbf{【小白补丁】#1}\par\nobreak\noindent\ignorespaces
}{%
  \par\noindent\rule{\linewidth}{0.4pt}\par\smallskip
}

\title{InstructGraph 中文精读导读：用“类代码”格式让大模型读懂图、少胡说}
\author{原文：Jianing Wang, Junda Wu, Yupeng Hou, Yao Liu, Ming Gao, Julian McAuley（ACL 2024 Findings）}
\date{}

\begin{document}
\maketitle
\begin{center}\small\textit{本导读为面向初学者的中文学习笔记：含要点翻译、公式精讲与图表解读。}\end{center}

\section*{论文卡片}
\begin{center}
\begin{tabular}{@{}ll@{}}
\toprule
\textbf{标题} & InstructGraph: Boosting Large Language Models via Graph-centric \\
            & Instruction Tuning and Preference Alignment \\
\textbf{作者} & Jianing Wang, Junda Wu, Yupeng Hou, Yao Liu, Ming Gao, Julian McAuley \\
\textbf{机构} & 华东师范大学、加州大学圣地亚哥分校（UC San Diego） \\
\textbf{会议} & ACL 2024 Findings（计算语言学协会年会·发现专辑） \\
\textbf{arXiv} & 2402.08785 \\
\textbf{代码} & \url{https://github.com/wjn1996/InstructGraph} \\
\textbf{方向} & 图（Graph）$\times$ 大模型 $\times$ 后训练（指令微调 + 偏好对齐） \\
\textbf{一句话} & 把所有图都写成统一的“类代码”格式喂给大模型，先指令微调学会图任务，\\
              & 再用 DPO 偏好对齐压幻觉，最终在图推理与图生成上超过 GPT-4。 \\
\bottomrule
\end{tabular}
\end{center}

\section{一、这篇论文在解决什么问题}

大语言模型（LLM，如 GPT、LLaMA）天生只会处理“一串文字”。但现实世界里很多重要信息天生是\textbf{图}：
知识图谱、符号图、社交网络，甚至人脑里隐式的“思维图”。这些图由\textbf{节点}（实体）和\textbf{边}（关系）
组成，很难用一段平铺直叙的文字说清楚。于是大家都想把“图”塞进 LLM。作者把已有做法归为两条路、并各指出一个痛点：

\begin{enumerate}[leftmargin=2em]
  \item \textbf{“图描述”路线}：人工写一大堆模板，把图“翻译”成自然语言句子喂给 LLM。
        痛点：模板要海量人工设计；更麻烦的是，LLM 生成的回答\textbf{很难再被解析回一张真正的图}。
  \item \textbf{“图嵌入”路线}：额外训练一个图神经网络（GNN），把图编码成向量，注入 LLM 的部分参数里。
        痛点：把一张大图或复杂图压成向量会\textbf{丢信息}，而且也难以把输出还原成图。
\end{enumerate}

在这两条路之外，作者还点出两个更本质的挑战：（1）\textbf{图与文本之间存在语义鸿沟}，
直接拼起来 LLM 未必能真正理解；（2）LLM 在图任务上\textbf{爱胡说（幻觉）}——
比如图里缺关键信息时硬答、或者生成一张事实错误／自相矛盾／信息缺失的图。作者把后者称为\textbf{图幻觉（graph hallucination）}问题。

\textbf{本文的目标}：提出一个叫 \textbf{InstructGraph} 的框架，分两步走——先用\textbf{图指令微调（graph instruction tuning）}
教会 LLM 做各种图任务，再用\textbf{图偏好对齐（graph preference alignment，DPO 风格）}压低图任务上的幻觉。
关键招数是：不外挂任何图编码器，而是用一种统一的\textbf{“类代码”结构化格式}把所有图写出来——
因为现在的 LLM 本来就很擅长读写代码。

\begin{tip}{为什么“把图写成代码”是个好主意}
LLM 在海量代码上预训练过，对 \texttt{list = [...]}、\texttt{(a -> b)} 这类结构化写法天然敏感。
把一张图写成 \texttt{entity\_list=[...]; triple\_list=[(...)];} 这种“类代码”，既保留了图的完整结构（不丢信息），
又落在 LLM 最擅长的格式里，输出还能被程序直接解析回一张图。一举解决了前面两条老路的通病。
\end{tip}

\section{二、摘要中译}

\begin{quote}
当前的大语言模型（LLM）通过参数更新能否更好地求解图推理与图生成任务？本文提出 \textbf{InstructGraph}，
一个通过\textbf{指令微调（instruction tuning）}与\textbf{偏好对齐（preference alignment）}赋予 LLM 图推理与图生成能力的框架。
具体而言，我们首先提出一种\textbf{结构化格式表述器（structured format verbalizer）}，把所有图数据统一成一种通用的
“类代码”格式，从而无需任何外部图专用编码器即可简洁地表示图。其次，我们引入\textbf{图指令微调}阶段，
引导 LLM 求解图推理与图生成任务。最后，我们识别出图任务中潜在的幻觉问题，并采样负样本进行偏好对齐，
目标是提升模型输出的可靠性。在多个以图为中心的任务上的大量实验表明，InstructGraph 能取得最佳性能，
分别超过 GPT-4 与 LLaMA2 达 13\% 和 38\% 以上。
\end{quote}

\section{三、核心贡献}
\begin{enumerate}[leftmargin=2em]
  \item \textbf{提出结构化格式表述器}：用统一的“类代码”格式表示任意图，\emph{不外挂图编码器}，
        既不丢结构信息，又对齐 LLM 最擅长的代码能力，输出可被解析回真正的图。
  \item \textbf{构建图指令微调语料与流程}：把 29 个图任务整理成约 160 万条指令数据，定义\textbf{四大类图任务}
        （图结构建模、图语言建模、图生成建模、图思维建模），统一用标准的因果语言建模目标继续训练 LLM。
  \item \textbf{提出图偏好对齐缓解幻觉}：针对图推理与图生成设计多种“图幻觉”场景（不实图、冲突图、缺失图等），
        自动负采样造“错误答案”，用 DPO 优化 LLM 提升可靠性。
  \item \textbf{大量实验验证}：在图指令任务、图偏好任务以及通用 NLP 任务上全面评测，
        超过同规模 LLaMA2/Vicuna，并在总体上超过 GPT-4。
\end{enumerate}

\section{四、预备知识（小白先看这里）}

\begin{tip}{什么是指令微调（instruction tuning）}
就是把“任务说明 + 输入 + 标准答案”整理成大量样本，继续训练 LLM，让它学会“听懂指令并照做”。
本文把各种图任务（判断连通、生成知识图谱、回答图问题……）都写成统一格式的指令样本，
用最普通的“预测下一个词”目标来训练，就能让一个通用 LLM 变成会做图任务的 InstructGraph-INS。
\end{tip}

\begin{tip}{什么是 DPO（直接偏好优化）与偏好对齐}
DPO（Direct Preference Optimization）是一种让模型“学会偏好好答案、嫌弃坏答案”的训练方法，
是 RLHF（基于人类反馈的强化学习）的轻量替代品：它\textbf{不需要单独训练奖励模型、也不跑强化学习}，
只要给一对“(好答案 $Y^+$, 坏答案 $Y^-$)”，直接用一个公式把 LLM 往“多给 $Y^+$、少给 $Y^-$”的方向推。
本文的妙处在于：好答案就是原本的正确答案，\textbf{坏答案靠“故意把图改错”自动生成}，不用人工标注。
\end{tip}

形式化记号（看不懂可跳过，用到时再回看）：设有 $M$ 个图任务 $D=\{D_1,\cdots,D_M\}$，
每个任务的数据集为 $D_j=\{(I_i,G_i,P_i,A_i)\}_{i=1}^{N_j}$，其中 $I_i$ 是\textbf{指令}，
$P_i$ 是可选的\textbf{文本段落}，$A_i$ 是\textbf{最终答案}；图 $G_i=(E_i,R_i,T_i,S_i)$ 由
节点（实体）集 $E_i$、可选关系集 $R_i$、边（三元组）集 $T_i$ 和可选文本属性集 $S_i$ 组成。
关键的“图$\to$文字”转换由\textbf{结构化格式表述器} $M(\cdot)$ 完成：
\begin{equation*}
  C_i = M(G_i). \tag{$\ast$}
\end{equation*}
也就是把整张图 $G_i$ 写成一段“类代码”序列 $C_i$，例如
\texttt{entity\_list=[...]; triple\_list=[(...)];}。带文本属性的节点用“面向对象”的写法表达，
例如把节点 \texttt{User1} 的评论写成 \texttt{User1.review=The film is nice.}。
这样所有图都被统一成同一种格式，可以和普通文字无缝拼接喂给 LLM。

\section{五、四大类图任务与数据集}

如图~\ref{fig:tasks-intro} 所示，作者把图相关任务归成\textbf{四大类}（前两类偏图推理，第三类是图生成，第四类兼有推理和生成）。
这套分类是全文的骨架，先记住它再看后面会顺很多。

\begin{table}[H]\centering
\caption{InstructGraph 的四大类图任务（整理自原文 Figure 1 与 Table 1）。}
\begin{tabular}{p{2.6cm}p{5.6cm}p{4.4cm}}
\toprule
\textbf{任务大类} & \textbf{包含的子任务} & \textbf{做什么} \\
\midrule
图结构建模\newline (Graph Structure Modeling) & 连通检测、环检测、哈密顿路径、二分图匹配、最短路径、度数计算 & 让 LLM 理解最基础的图结构（只给节点、有/无向边、可选权重）。\emph{纯图推理} \\
\midrule
图语言建模\newline (Graph Language Modeling) & 图标题生成、图问答、节点分类、链接预测、相关性检测、协同过滤 & 让 LLM 同时理解图的\emph{结构}与\emph{语义}并回答问题。\emph{纯图推理} \\
\midrule
图生成建模\newline (Graph Generation Modeling) & 知识图谱生成、结构图生成 & 给一段文字，让 LLM \emph{生成}一张“类代码”格式的图。\emph{图生成} \\
\midrule
图思维建模\newline (Graph Thought Modeling) & 算术、符号、机器人、逻辑推理 & 三步走：①找问题主体 ②生成一张“思维图”表达推理过程 ③据图给出答案。\emph{推理+生成，仅用于评测} \\
\bottomrule
\end{tabular}
\end{table}

\begin{figure}[H]\centering
  \includegraphics[width=0.82\linewidth]{cand_p02_k00.pdf}
  \caption{原文 Figure 1：四大类以图为中心的推理与生成任务示例（图内英文为原文保留）。
  左上为图结构建模（连通/环/最短路等），右上为图语言建模（图问答/节点分类/链接预测等），
  左下为图生成建模（知识图谱/结构图生成），右下为图思维建模（把推理过程画成思维图再作答）。
  注意每张图都被写成 \texttt{node\_list}/\texttt{edge\_list}/\texttt{triple\_list} 这样的“类代码”格式。}
  \label{fig:tasks-intro}
\end{figure}

\textbf{数据规模}：作者一共收集/构造了 \textbf{29 个图任务}，组成约 \textbf{160 万条}指令微调数据
和约 \textbf{10 万条}偏好对齐数据（精确统计见原文 Table 10：指令训练集约 134 万、偏好训练集约 8.6 万）。
数据来源很广，覆盖 NLGraph（图结构）、WebNLG/GenWiki（图标题）、WebQSP/GrailQA（图问答）、
Cora/Citeseer/Pubmed/Arxiv/Products（节点分类）、Wikidata/FB15K-237/ConceptNet（链接预测）、
UIE（信息抽取/知识图谱生成）等经典数据集。

\begin{tip}{“图思维建模”为什么只用来评测、不参与训练}
作者只用前三类任务做指令微调；第四类“图思维建模”（让模型把算术/逻辑题的推理过程画成一张思维图）
\textbf{故意不放进训练集}，专门留作测试。这样可以检验：在前三类图任务上学到的能力，
能不能\emph{迁移}去帮助原本没训练过的推理任务——这是对“图能力是否真正泛化”的一次考验。
\end{tip}

\section{六、方法逐节精讲}

InstructGraph 的整体骨架（图~\ref{fig:arch}）由\textbf{三个模块}串成一条流水线：
① 图输入工程（Graph Input Engineering）→ ② 图指令微调（Graph Instruction Tuning）→ ③ 图偏好对齐（Graph Preference Aligning）。

\begin{figure}[H]\centering
  \includegraphics[width=0.98\linewidth]{cand_p03_k00.pdf}
  \caption{InstructGraph 总体框架（原文 Figure 2，图内英文为原文保留）：
  \textbf{左下「Graph Input Engineering」}——用结构化格式表述器（\texttt{</> Verbalizer}）把图写成“类代码”格式；
  \textbf{上方「Graph Instruction Tuning」}——把约 160 万条图任务数据按推理/生成喂给 LLM，
  用最大似然（Maximum Likelihood）做指令微调，仅训练 LoRA 适配器（图中 $\heartsuit$ 表示可训练、$\ast$ 表示冻结主干）；
  \textbf{右下「Graph Preference Aligning」}——构造四种“图幻觉”场景的错误答案作为负样本，
  用 DPO 优化算法把模型推向更可靠的输出。}
  \label{fig:arch}
\end{figure}

\subsection{第 1 步：图输入工程 —— 结构化格式表述器}
这是本文第一块技术核心，回答“怎么把图对齐到文本、塞进 LLM 的序列接口”。
做法已在第四节讲过：用表述器 $M(\cdot)$ 把图 $G_i$ 映射成“类代码”序列 $C_i=M(G_i)$。
基本格式是：所有节点（或实体）列进 \texttt{node\_list}（或 \texttt{entity\_list}），
所有边（或三元组）列进 \texttt{edge\_list}（或 \texttt{triple\_list}）；
带附加信息的节点用“面向对象”写法表达，如 \texttt{User1.review=The film is nice.}。
\textbf{关键好处}：这种格式不丢图的任何结构，又能被 LLM 轻松读写、被程序解析回真图。

\subsection{第 2 步：图指令微调（得到 InstructGraph-INS）}
经过图输入工程后，每条样本都变成了纯文本序列，于是可以\textbf{直接复用最标准的因果语言建模（CLM）目标}
继续训练 LLM——不需要任何特殊网络结构。对任务数据集 $D_j=\{(I_i,G_i,P_i,A_i)\}_{i=1}^{N_j}$，
用\textbf{最大似然}优化，让模型在给定输入 $X_i$ 时尽量输出参考答案 $A_i$：
\begin{equation*}
  \mathcal{L}(D_j) = -\sum_{i=1}^{N_j} \log \pi_\theta\big(Y_i = A_i \mid X_i\big). \tag{1}
\end{equation*}
其中 $\pi_\theta$ 是参数为 $\theta$ 的 LLM，$Y_i$ 是模型输出，$X_i$ 是输入序列，$A_i$ 是参考答案。
不同任务的输入/输出拼法不同（见原文 Table 1）：比如图结构任务 $X_i=[I_i,C_i]$、$Y_i=A_i$；
图生成任务 $X_i=[I_i,P_i]$、\textbf{$Y_i=C_i$}（即让模型\emph{生成}那段类代码图）。训练完得到 \textbf{InstructGraph-INS}。

\begin{tip}{怎么读公式 (1)：它其实就是“背标准答案”}
$\log\pi_\theta(Y_i=A_i\mid X_i)$ 表示“在输入 $X_i$ 下，模型给出正确答案 $A_i$ 的对数概率”。
前面加负号、再求和，就是让这些概率尽量大（损失尽量小）——也就是逼着模型在每个图任务上都尽量输出对的答案。
这是所有指令微调最朴素的写法，本文的创新不在这个公式，而在于\emph{喂进去的 $X_i$ 是统一的“类代码”图格式}。
\end{tip}

\subsection{第 3 步：图偏好对齐 —— 用 DPO 压幻觉（得到 InstructGraph-PRE）}
这是本文第二块技术核心。光会做题还不够，模型有时会\textbf{幻觉}：图里缺信息时硬编答案、或生成一张事实错误的图。
作者为图推理和图生成分别设计了\textbf{多种幻觉场景}，并\textbf{自动负采样}造出“坏答案 $Y_i^-$”：

\begin{itemize}[leftmargin=2em]
  \item \textbf{图推理的四种幻觉}：① 图正确但答错；② 不实图（unfactual，与外部知识矛盾）导致答错；
        ③ 冲突图（conflict，图内信息自相矛盾）导致答错；④ 缺失图（missing，缺关键信息）导致答错。
  \item \textbf{造负样本的手法}：对场景①，从别的样本里随机抓一个答案当 $Y_i^-$；
        对其余场景，随机\textbf{替换/增加/删除}图里的一些节点或边，构造出“被改坏的图”。
        此时\textbf{原来的正确答案反而成了负样本 $Y_i^-$}，而正样本 $Y_i^+$ 是一句话：
        “抱歉，输入图含有错误信息，所以该问题无法直接回答。”——即教模型学会“图不对就拒答/纠错”。
  \item \textbf{图生成的幻觉}：图生成更难（要输出一整段完整准确的类代码）。同样用替换/增删操作直接造出错误图当 $Y_i^-$，
        原始正确图作为 $Y_i^+$。
\end{itemize}

有了成对的 $(X_i, Y_i^+)$ 和 $(X_i, Y_i^-)$，先按 \textbf{Bradley-Terry} 偏好模型定义“偏好概率”：
\begin{equation*}
\begin{aligned}
  p_\theta(Y_i^+ > Y_i^- \mid X_i) &= \frac{1}{1+\exp\{\,r(Y_i^+,Y_i^-,X_i)\,\}}, \\
  r(Y_i^+,Y_i^-,X_i) &= -\,\beta\log\frac{\pi_\theta(Y_i^+\mid X_i)}{\pi_{\mathrm{ref}}(Y_i^+\mid X_i)}
  + \beta\log\frac{\pi_\theta(Y_i^-\mid X_i)}{\pi_{\mathrm{ref}}(Y_i^-\mid X_i)}.
\end{aligned}\tag{2}
\end{equation*}
其中 $\beta$ 是平衡因子，$\pi_\theta$ 是待训练的\textbf{策略模型}、$\pi_{\mathrm{ref}}$ 是固定的\textbf{参考模型}
（都用上一步的 InstructGraph-INS 初始化）。最终用最大似然优化 DPO 目标：
\begin{equation*}
  \mathcal{J}(\pi_\theta,\pi_{\mathrm{ref}}) = -\,\mathbb{E}_{(X_i,Y_i^+,Y_i^-)\sim D}
  \left[\log\sigma\!\left(\beta\log\frac{\pi_\theta(Y_i^+\mid X_i)}{\pi_{\mathrm{ref}}(Y_i^+\mid X_i)}
  - \beta\log\frac{\pi_\theta(Y_i^-\mid X_i)}{\pi_{\mathrm{ref}}(Y_i^-\mid X_i)}\right)\right]. \tag{3}
\end{equation*}
训练得到的策略 $\pi_\theta$ 即 \textbf{InstructGraph-PRE}。

\begin{tip}{怎么读 DPO 目标 (3)：一句话拆给你听}
看括号里那两项之差：$\beta\log\frac{\pi_\theta(Y^+)}{\pi_{\mathrm{ref}}(Y^+)}$ 衡量“相比初始模型，现在更愿意给好答案吗”，
后一项是“更愿意给坏答案吗”。两者相减再过 $\sigma$（Sigmoid）取对数：当模型\textbf{把好答案概率提上去、坏答案概率压下去}时，
括号里为正、$\log\sigma$ 接近 0（损失小）；反之损失大。所以最小化 $\mathcal{J}$ 就等价于
“相对参考模型，多给 $Y^+$、少给 $Y^-$”。$\pi_{\mathrm{ref}}$ 像一根锚，防止模型为了讨好偏好而跑偏、忘掉原有能力。
而 $\beta$ 控制“可以偏离锚多远”。这正是 DPO 不需要奖励模型、不跑强化学习就能做偏好对齐的精髓。
\end{tip}

\textbf{训练细节}：默认主干为 LLaMA2-7B-HF，最大长度 2048，优化器 AdamW；
指令微调阶段学习率 $5\mathrm{e}{-5}$、偏好对齐阶段降到 $5\mathrm{e}{-7}$。
为省显存用了 FSDP + CPU Offloading + FlashAttention + BFloat16，
并用 \textbf{LoRA}（rank$=32$，lora\_$\alpha=128$）做参数高效微调——\textbf{只训练很少的参数，冻结 LLM 主干}。

\section{七、实验与结果解读}

\textbf{对比对象}：同规模开源模型 LLaMA2-7B、Vicuna-7B；以及强基线 GPT-3.5 (turbo) 和 GPT-4。
评测主要在\textbf{零样本（zero-shot）}下进行。

\subsection{主结果一：图指令任务（超过 GPT-4）}
原文 Table 2 在 29 个图任务上对比，InstructGraph-INS 的\textbf{总平均分 79.84\%}，
比 GPT-4 的 66.76\% 高出 \textbf{13.08\%}，比同规模 LLaMA2（41.65\%）高出约 38\%。下面节选几行最有代表性的数字：

\begin{table}[H]\centering
\caption{图指令任务主结果（\%，数值源自原文 Table 2，越高越好；本表为节选+总平均，加粗为该行最优）。}
\begin{tabular}{llccccc}
\toprule
任务 & 指标 & GPT-3.5 & GPT-4 & LLaMA2 & Vicuna & \textbf{InstructGraph-INS} \\
\midrule
环检测 Cycle Dect. & ACC & 59.02 & 61.44 & 50.79 & 52.88 & \textbf{91.10} \\
二分图匹配 Bipt. Match & ACC & 50.23 & 66.11 & 0.00 & 0.00 & \textbf{76.36} \\
最短路径 Shrt. Path & ACC & 38.99 & 49.03 & 0.00 & 0.00 & \textbf{66.29} \\
图问答 PathQSP & EM & 52.54 & 68.64 & 42.70 & 31.90 & \textbf{86.40} \\
图问答 WebQSP & EM & 53.73 & 61.57 & 40.07 & 26.42 & \textbf{73.30} \\
链接预测 Wikidata & Hits@1 & 43.73 & 62.94 & 10.75 & 10.38 & \textbf{96.52} \\
链接预测 FB15K-237 & Hits@1 & 60.34 & 66.88 & 0.00 & 0.00 & \textbf{98.91} \\
推荐 Amazon & Hits@1 & 27.09 & 59.77 & 44.40 & 16.40 & \textbf{78.80} \\
信息抽取 UIE & F1 & 24.41 & 26.22 & 20.21 & 26.11 & \textbf{76.82} \\
\midrule
\textbf{总平均 Avg.} & & 59.45 & 66.76 & 41.65 & 46.06 & \textbf{79.84} \\
\bottomrule
\end{tabular}
\end{table}

\textbf{怎么读这张表}：（1）在“图结构”类任务上，LLaMA2/Vicuna 经常考\textbf{0 分}（完全不会做二分图匹配、最短路），
而 InstructGraph 普遍 60\%$\sim$90\%；说明这些结构推理能力确实是被\emph{专门训练}出来的。
（2）在链接预测上 InstructGraph 高达 96\%$\sim$99\%，远超 GPT-4，体现“类代码三元组格式”对关系预测的天然契合。
（3）作者也坦诚：在 Degree Computing、WebNLG、GenWiki、WikiTQ、Citeseer 等少数任务上 InstructGraph 略逊于 GPT-3.5/GPT-4，
因为超大参数模型“记”住了更多相似知识；但在其余推理任务上 InstructGraph 仍领先约 10\%。

\subsection{主结果二：图偏好任务（DPO 进一步压幻觉）}
作者构造了一个“幻觉测试集”：每条样本含一个对答案和一个错答案，计算 LLM 在两者上的\textbf{困惑度（PPL）}，
选 PPL 更低（更“顺”）的那个作为模型的偏好结果，于是\textbf{准确率越高=越能识别正确答案=幻觉越少}。

\begin{table}[H]\centering
\caption{图偏好任务主结果（\%，数值源自原文 Table 3，越高越好/越忠实；加粗为最优）。}
\begin{tabular}{lcccccc}
\toprule
方法 (7B) & 是否对齐 & 结构 & 标题 & 图问答 & 节点分类 & 信息抽取 \\
\midrule
LLaMA2 & 否 & 38.64 & 57.96 & 70.70 & 74.68 & 37.40 \\
Vicuna & 否 & 39.12 & 62.37 & 64.38 & 77.63 & 40.80 \\
InstructGraph-INS & 否 & 50.32 & 81.15 & 77.85 & 83.16 & 69.14 \\
\textbf{InstructGraph-PRE} & \textbf{是} & \textbf{57.80} & \textbf{87.44} & \textbf{84.44} & \textbf{88.98} & \textbf{91.44} \\
\bottomrule
\end{tabular}
\end{table}

\textbf{解读}：仅做指令微调的 InstructGraph-INS（平均 72.32\%）就已经比 LLaMA2/Vicuna 高约 16\%、15\%，
说明\textbf{把任务知识注入参数本身就能减少幻觉}。在此之上再叠加 DPO 偏好对齐，
InstructGraph-PRE 平均再涨约 10\%（达 82.02\%），尤其在信息抽取上从 69.14 飙到 91.44——
证明\textbf{精心设计的偏好优化能进一步逼近上限、教会模型“图不对就别硬答”}。

\subsection{其他关键发现}
\begin{itemize}[leftmargin=2em]
  \item \textbf{“类代码”格式 vs 传统模板}（原文 Table 7）：把图写成结构化代码，全面优于 InstructGLM 式的自然语言模板；
        \textbf{且传统模板根本无法支持图生成}，而类代码格式可以。这直接验证了核心设计的价值。
  \item \textbf{消融实验}（原文 Table 8）：去掉任一类建模任务都会掉分——图语言建模对“图问答/节点分类”最关键，
        图生成建模对“信息抽取”最关键；在偏好侧，\textbf{去掉“缺失图”负样本掉分最多}，
        说明\emph{图里缺关键信息是幻觉的最大来源}。
  \item \textbf{参数高效微调}（原文 Figure 4）：在六种参数规模下对比 LoRA、Prefix-tuning、Adapter、全量微调，
        \textbf{LoRA 用极少可训练参数就逼近全量微调}，性价比最高。
  \item \textbf{换骨干/换规模}（原文 Figure 5）：在 LLaMA2-7B/13B、Vicuna-7B/13B 上 InstructGraph 都能带来大幅提升，
        且更大模型收益更大。有趣的是：微调前 Vicuna 略强于 LLaMA2，但\textbf{图指令微调后这个优势被反转}，
        作者分析 LLaMA2 是更好的“起点”骨干。
  \item \textbf{思维规划迁移}（原文 Table 5）：在算术/符号/机器人/逻辑等\emph{未训练}的推理任务上，
        用 CoT 或“图思维建模（GTM）”提示，InstructGraph-INS 都优于 LLaMA2/Vicuna，
        说明图能力能迁移去帮助通用推理。
  \item \textbf{不伤通用能力}（原文 Table 4、Table 6）：在 HaluEval、TruthfulQA、BBH、MMLU 等通用 NLP 基准上，
        InstructGraph 仍与同规模开源模型持平或略优，说明\textbf{专攻图任务没有把通用语言能力练废}。
\end{itemize}

\section{八、小白知识点卡片（速查）}
\begin{itemize}[leftmargin=2em]
  \item \textbf{结构化格式表述器（structured format verbalizer）}：本文核心，把任意图写成统一“类代码”
        （\texttt{entity\_list}/\texttt{triple\_list}）的转换器，\emph{不外挂图编码器}。
  \item \textbf{指令微调（instruction tuning）}：用“指令+输入+答案”样本继续训练 LLM，让它学会按指令做任务。
  \item \textbf{DPO（直接偏好优化）}：给一对“好/坏答案”，不训练奖励模型、不跑强化学习，直接把模型推向好答案的偏好对齐方法。
  \item \textbf{Bradley-Terry 模型}：把“A 比 B 更好”的偏好建模成一个概率的经典统计模型，是 DPO 的理论基础。
  \item \textbf{图幻觉（graph hallucination）}：模型在图任务上编造不存在的节点/边，或图缺信息时硬答。
  \item \textbf{负采样（negative sampling）}：本文靠“随机替换/增删图中节点或边”自动造出错误答案当负样本，省去人工标注。
  \item \textbf{LoRA / PEFT}：参数高效微调；LoRA 只训练低秩增量矩阵，冻结主干、省显存。
  \item \textbf{因果语言建模（CLM）}：最标准的“预测下一个词”训练目标。
  \item \textbf{Hits@1 / EM / PPL}：分别为“排第一是否命中”“完全匹配率”“困惑度（越低越顺）”。
\end{itemize}

\section{九、一句话总结与“图 $\times$ 后训练”谱系定位}

\textbf{一句话}：InstructGraph 不给 LLM 外挂任何图编码器，而是用一种统一的\textbf{“类代码”结构化格式}把所有图写出来，
先做\textbf{图指令微调}让模型学会 29 个图任务，再用\textbf{DPO 偏好对齐}（靠“故意改错图”自动造负样本）压低图幻觉，
最终在图推理与图生成上\emph{超过 GPT-4 与 LLaMA2}。

\textbf{在“图 $\times$ 后训练”谱系里的位置}：本系列各篇的关键区别，就在于\emph{“把图的信息塞进 LLM 的方式”}不同——
G-Retriever 走“检索 + 软提示（PCST 抠子图）”，LLaGA 走“投影器对齐（把图向量投影进词空间）”，
G1 走“强化学习”。而 \textbf{InstructGraph 走的是“纯文本化（类代码 verbalizer）+ 指令微调 + DPO 偏好对齐”路线}：
它\textbf{不引入任何图神经网络}，把图问题彻底转化成 LLM 最擅长的“读写代码 + 跟随指令”问题，
再用偏好对齐这一步专门治图任务的幻觉。可以把它和 GraphWiz 放在一起看——两者都属于
“\textbf{指令微调 + 偏好/DPO}”这一支，代表了“不改 LLM 结构、只靠数据与后训练把图能力灌进去”的思路。

\textbf{动手复现}：官方代码 \url{https://github.com/wjn1996/InstructGraph}
（含结构化格式表述器、29 个图任务的语料构造、指令微调与 DPO 偏好对齐流程），
是想理解“图 $\times$ 指令微调 $\times$ DPO”这条路线很好的练手项目。

\end{document}
